% =============================================================
%  Supplementary Material for
%  "Reading the Banished: A Deep-Learning Emotion Analysis of
%   Tang--Song Exiled Literati"
%  Repository: https://github.com/fourhairsall/tangsong-exile-affect
%  Compile: xelatex tangsong_exile_affect_supplement.tex
% =============================================================
\documentclass[11pt]{article}
\usepackage[UTF8]{ctex}
\usepackage{geometry}
\geometry{a4paper,margin=1in}
\setlength{\emergencystretch}{2em}
\usepackage{booktabs}
\usepackage{multirow}
\usepackage{array}
\usepackage{amsmath}
\usepackage{amssymb}
\usepackage{graphicx}
\usepackage{url}
\usepackage{xcolor}
\usepackage{float}
\usepackage{natbib}
\setlength{\bibsep}{0.5ex}
\bibliographystyle{plainnat}
\renewcommand{\figurename}{Figure}
\renewcommand{\tablename}{Table}
\renewcommand{\refname}{References}

\title{Supplementary Material\\[0.4em]
\large Reading the Banished: A Deep-Learning Emotion Analysis of
Tang--Song Exiled Literati}
\author{Shujiang Tang\\
Hunan University of Science and Engineering\\
\texttt{sjtang@huse.edu.cn}}
\date{}

\begin{document}
\maketitle

\noindent This supplement carries the material referred to from the main text as
\emph{Supplementary Material}: the corpus-construction record (App.~A), the
cross-edition invariant-representation probe (App.~B), and the lexicon, exile
chronology and gold-annotation protocol (App.~C). It is self-contained and cites the
main paper only by name.

\section{Corpus Construction: Transcription, Parsing and Attribution}
\label{app:parse}
\emph{Yu Xuan Song Shi} survives only as the mandoku plaintext of \emph{Yu Xuan Song
Jin Yuan Ming Si Chao Shi} (kanripo KR4h0143), in which the 78 Song \emph{juan}
interleave author registers, prefaces and a handful of Jin/Ming \emph{juan}. We
parsed every \emph{shi} body line (the lowest-indent text) and attributed it to its
enclosing author by exact-matching the (simplified) author name against the
\emph{Quan Song Shi} author registry of 8{,}925 names; this yields 9{,}405 poems by
660 authors, and \emph{all} 660 match an author present in QSS, confirming the
parse. \emph{Song Shi Chao} is transcribed the same way: verse is the lowest-indent
text, each poet's section is anchored on the \emph{ji}-name heading---including the
上/中/下 continuation headings that carry 欧阳修, 王安石 and 苏轼 across several
volumes---and the prose biographies that habitually follow those headings are
discarded rather than counted as verse, since they narrate demotion and transfer and
would inflate the 贬谪羁旅 and 悲苦孤寂 registers on the anthology side alone. The
parse anchors 98 poet headings and yields 9{,}612 poems by 76 authors over 829{,}341
characters of verse, with no unresolved attribution.

Two documented features of this text bear on how the comparison should be read. The
\emph{Siku} digest records that sixteen of the poets it lists appear here as headings
without text (``有録無書\ldots 蓋剞劂未竣故竟無完帙也''), and describes the anthology
as gathering 百家; if exact, this state should yield eighty-four poets with verse,
and our 76 leave a residue of eight that we neither anchor nor explain away. The
shortfall affects author \emph{coverage} only: $\delta$ is computed over authors
present in both sources and is not directly biased by it. The same digest also reports
that because each poet's section was printed as a separate booklet, copies in
circulation differ in extent (``所傳之本往往多寡不同'')---the paper's own thesis
stated by a Qing editor, for ``the edition'' of a selection anthology is not a stable
object even before measurement begins.

\label{app:attr}
Attribution---identifying \emph{who} wrote a poem---is a separate question from whether
the affect measurement is trustworthy, and we keep it out of the main line. The task is
well studied: Cap-Transformer models for Tang-poet identification \citep{shiziren2022},
LDA-Transformer mixtures \citep{zhou2021lda}, and imagery-based style separating 柳永
from 苏轼 \citep{fang2011cld}. We compare three classifiers on identical 85/15
stratified splits under one fixed seed: a character CNN \citep{kim2014convolutional},
character $n$-gram TF-IDF with logistic regression, and the same features with a linear
SVM. Table~\ref{tab:attr} shows the lightweight TF-IDF + SVM \emph{beats} the CNN on
every task and is more reproducible and GPU-free; the 12-way author task is held back
by severe class imbalance (minor authors 154--320 poems, major authors 2{,}000+), which
gives low macro-F1 despite high weighted-F1.

\begin{table}[H]
\centering\footnotesize
\caption{Attribution: accuracy / macro-F1. Best per row in bold.}
\label{tab:attr}
\begin{tabular}{lccc}
\toprule
Task & CNN & TF-IDF+LR & TF-IDF+SVM\\
\midrule
Author (12) & .624/.271 & .714/.351 & \textbf{.802/.568}\\
Dynasty (2) & .894/.837 & .895/.808 & \textbf{.936/.898}\\
Tang prose (4) & .619/.512 & .856/.812 & \textbf{.924/.908}\\
\bottomrule
\end{tabular}
\end{table}

\section{Cross-Edition Invariant-Representation Probe (M6)}
\label{app:m6}
As a probe---not a core claim---we ask whether a \emph{learned} representation could
hold a poet's affect while shedding edition-specific leakage. We freeze the PPMI--SVD
character embeddings (100-d, unit-norm), represent each poem by the mean-pool of its
character vectors, and take the author-mean vector for every author present in
\emph{both} editions with $\ge$300 characters. Affect is read out as in
the transmission-robustness analysis of the main paper: project each author vector onto the five lexicon pole directions
and compute $(A,S)$. We obtain a single linear ``source-invariant'' projector by a
\emph{closed-form domain-adversarial} solution---train a logistic adversary to predict
edition from the author-mean vector, then project out its weight direction,
$g(\mathbf x)=\mathbf x-(\mathbf x\!\cdot\!\mathbf w)\mathbf w$---and apply it to both
settings.

\begin{table}[H]
\centering\footnotesize
\caption{M6 invariant-representation probe. Frozen = identity projector on the
PPMI--SVD author-mean backbone; M6 = closed-form domain-adversarial projector.
$A$/$S$ are the concordance coefficients read out from the 5-dim pole projection.
edAUROC = 5-fold logistic edition predictability (lower is more invariant); varR =
dimensionwise variance retained by M6 ($1.0$ = none removed).}
\label{tab:m6}
\setlength{\tabcolsep}{3.2pt}
\begin{tabular}{lcccccc}
\toprule
Setting & $n$ & Variant & $A$ & $S$ & edAUROC & varR\\
\midrule
Tang: QTS vs.\ YD & 363 & frozen & $0.865$ & $1.00$ & $0.536$ & $1.00$\\
Tang: QTS vs.\ YD & 363 & M6 & $0.981$ & $1.00$ & $0.537$ & $0.69$\\
\midrule
Song: QSS vs.\ YXSS & 216 & frozen & $0.900$ & $1.00$ & $0.848$ & $1.00$\\
Song: QSS vs.\ YXSS & 216 & M6 & $0.980$ & $1.00$ & $0.584$ & $0.49$\\
\bottomrule
\end{tabular}
\end{table}

On the Tang same-work pair the edition direction is negligible to begin with (edAUROC
$0.536$ $\approx$ random), so M6 has almost nothing to remove. On the Song selection
pair it is large and dominant (frozen edAUROC $0.848$); M6 removes precisely that axis,
collapsing edition predictability to $0.584$ while retaining all five affect dimensions
($S{=}1.00$) and raising $A$ from $0.900$ to $0.980$. This is a \emph{probe}: its
$A$/$S$ values are \emph{not} directly comparable to the cross-edition concordance table of the main paper, and
its only licensed reading is qualitative---edition information here is a single,
low-rank, separable direction removable without harming the affect dimensions,
consistent with the main finding that edition effects act along specific directions
rather than as diffuse noise.

\section{Lexicon, Chronology and Annotation Protocol}
\label{app:lexicon}
The lexicon $\mathcal{L}$ partitions words into five categories; the \emph{full}
lists are recorded below. Each list mixes single characters with multi-character
words; all entries are stored in simplified script, consistent with the corpora's
normalisation (traditional$\to$simplified via OpenCC). The two affective poles
are marked (A); the thematic
exile--journey category (A$^\dagger$,
the named condition rather than a felt affect) is marked separately; the two imagistic
registers are marked (I).

\smallskip
\noindent\textbf{悲苦孤寂} \textsc{Bei} (A; sorrow--loneliness):
悲 伤 愁 怨 苦 孤 寂 独 凄 惨 哀 痛 恨 忧 戚 惘 怆 涕 泪 殇 殁 愍 悴 惸 黯 销 魂 零 萧 雁 砧 捣 吏 灯 店 舫 猿 驿 帆 仆 僮
断肠 凄凉 寂寥 孤危 憔悴 销魂 悲秋 苦辛 酸辛 辛酸 飘零 飘蓬 断蓬 飞蓬 转蓬

\smallskip
\noindent\textbf{贬谪羁旅} \textsc{Bian} (A$^\dagger$; exile--journey, the named condition rather than a felt affect):
贬 谪 迁 逐 流 羁 旅 客 宦 远 荒 蛮 瘴 殊 遐 徼 陬 裔
去国 天涯 海角 殊方 投荒 远谪 迁客 楚客 南荒 炎荒 瘴江 蛮烟

\smallskip
\noindent\textbf{旷达闲适} \textsc{Kuang} (A; detached--easeful):
闲 适 达 旷 豁 放 醉 啸 傲 悠 淡 泊 乐 欣 悦 遣 兀 疏 狂 从 容 恬 静 禅 苔 钓 隐 亭 栽 莲 畦 眠
旷达 闲适 放达 自适 陶然 悠然 忘机 委运 随遇 达观 适意

\smallskip
\noindent\textbf{自然山水} \textsc{Zi} (I; nature--landscape):
山 水 风 月 林 泉 石 竹 松 江 河 湖 海 花 鸟 溪 岩 岚 汀 洲 屿 峰 崖 涧 泽 涯 渚 苹 蓼
清风 明月 白云 沧浪 烟霞 苇岸 渔樵 桃源 落英 幽篁 暮云 孤云 寒云 春云 秋云 停云 烟云 云霞 云峰 云山

\smallskip
\noindent\textbf{空幻时空} \textsc{Kong} (I; void--temporality):
空 幻 梦 影 尘 虚 浮 瞬 逝 昔 古 千 秋 万 须 臾 泡 电 光 隙 驹 阴
浮生 如梦 梦幻 须臾 万古 古今 今古 泡影 微尘 刹那 弹指 过眼 逆旅 倏忽

\smallskip
\noindent\textbf{Corpus coverage of the lexicon.}
Table~\ref{tab:lexfreq} reports how often each category fires on the unified
\emph{shi} corpus (21{,}287 poems, 1{,}466{,}827 characters). All five categories
fire frequently and in reasonable balance; 自然山水 is expectedly the most
frequent, reflecting the centrality of
landscape in classical verse.

\begin{table}[H]
\centering\footnotesize
\caption{Lexicon corpus coverage (unified \emph{shi} layer).}
\label{tab:lexfreq}
\begin{tabular}{lrrrr}
\toprule
Category & Entries & Hits & Poems hit (\%) & Rate / 1k\\
\midrule
悲苦孤寂 Bei & 56 & 18,677 & 10,303 (48.4\%) & 12.73\\
贬谪羁旅 Bian & 30 & 9,549 & 6,721 (31.6\%) & 6.51\\
旷达闲适 Kuang & 43 & 19,446 & 11,054 (51.9\%) & 13.26\\
自然山水 Zi & 49 & 45,254 & 16,478 (77.4\%) & 30.85\\
空幻时空 Kong & 36 & 26,442 & 13,308 (62.5\%) & 18.03\\
\bottomrule
\end{tabular}
\end{table}

\noindent The ABI is
$\mathrm{ABI}(a)=\phi_{a,\textsc{Kuang}}-\phi_{a,\textsc{Bei}}$
(Eq.~2).

\begin{sloppypar}
\noindent The quietist (静禅) register is carried by the 旷达 detached--easeful
class via 禅/苔/钓/隐/亭/栽/畦/眠; the exile--journey and homesickness registers
are carried by the 悲苦 sorrow--loneliness class via 驿/帆/舫/店/灯/砧/捣/仆/僮
(entry counts in \texttt{results\_lexicon\_v2\_curation.json}: 悲苦孤寂
37$\to$47$\to$56, 旷达闲适 33$\to$43).
\end{sloppypar}

\label{app:chron}
None of the four corpora carries per-poem dating metadata (0 of 356{,}714 poems
across QTS/QSS/YD/YXSS carries a usable year field; the \emph{tags} of the Tang
files are thematic labels and the \emph{strains} of the Song files are tonal
patterns), so within-author period assignment required an external chronology
layer. For each queue official a set of exile events was compiled from the
standard scholarship (biographies, chronologies, and annotated editions), each
graded A (histories plus a dated chronology agree), B (scholarly consensus), or C
(disputed; never merged into A/B). Poems were then attached to periods by two
routes: internal chronological ordering, verified against anchor poems whose
dates are certain (for Su Shi the \emph{Quan Song Shi} sequence is chronological
from its first poem; the Crow-Terrace poem, \emph{First Arriving in Huangzhou},
\emph{Leaving Huangzhou}, \emph{First Arriving in Huizhou}, the Hainan poems and
the last poem all fall in strict order; the supplementary block after the final
poem is not chronological and was marked \emph{unknown}), or poem-level
attachment from dated annotated editions where the collection is ordered by
genre rather than by year. A $\pm1$-year buffer around each event boundary is
excluded from all contrasts. Coverage: Su Shi 473 exile / 995 pre / 860 post
poems (85\% of his corpus attached at grade A); Bai Juyi 59 / 66 / 18 at grade B
with a documented genre-selection caveat; the other four officials fell below
pairing power on attachment (Liu Yuxi 42 exile poems with no separable
pre/post control; Ouyang Xiu 12; Wang Yucheng 29; Huang Tingjian 4). Attachments, grades and coverage are released in
the replication package (\texttt{poem\_period\_labels.csv},
\texttt{exile\_chronology.csv}).

\label{app:goldrep}

\paragraph{Sampling.}
The gold sample is 160 \emph{shi} poems, 80 per dynasty, drawn by stratified random
sampling (seed 20260917) over the cross of dynasty, membership in the twelve-poet
cohort, and whether the poem's lexicon ABI falls above or below the median; within
each dynasty 60 poems come from the complete collection and 20 from a selection
anthology, so that the sample spans both transmission regimes and the full range of
the instrument's output rather than only its middle. Poems are restricted to 20--160
characters; below 20 characters a per-1{,}000-character rate is meaningless, and
above 160 characters the burden of close reading per item makes a blind double
coding impractical. The realised sample covers 78 authors. Coders received the poems
in a fixed random order with author, title, dynasty, anthology membership and all
lexicon output withheld, so the coding cannot be anchored on prior reputation.

\paragraph{Instrument and reliability.}
Each poem is rated on the five categories on a 0--3 ordinal scale (0 = absent, 1 =
present once incidentally, 2 = clearly present, 3 = dominant), with the category
definitions of \S\ref{app:lexicon} and two anchor examples per category agreed before
coding began; both coders rated all 160 poems independently and no disagreement was
resolved by discussion, so the reliability figures are not inflated by consensus.
Agreement is reported as quadratic-weighted $\kappa$, ordinal Krippendorff $\alpha$
and ICC$(2,1)$, and the gold index as the coder mean, exactly as in Eq.~(2).

\paragraph{Why split-half here is contiguous, not odd--even.}
Per-poem reliability of the lexicon is computed by applying the lexicon separately to
the first and second half of each poem's characters and reporting the Spearman--Brown
correction of the resulting correlation. Splitting by alternating characters would
sever every bigram and every parallel pair (对仗), driving the correlation to zero for
reasons that have nothing to do with the validity of a word-counting instrument; the
estimate is therefore conservative in construction and its near-zero value should be
read as a statement about the sparsity of poem-level word counts, not about the
lexicon's design. The same caution applies to $|\widehat{\rho}|$: a lexicon that
fires a handful of times in a 40-character poem cannot be reliable at that unit
whatever its aggregate behaviour.

\paragraph{Gold-set sensitivity and per-author agreement.}
Dropping poems shorter than 28 characters ($n{=}150$), dropping the extreme 5\% of
lexicon scores ($n{=}147$) or dropping the longest poems ($n{=}150$) leaves the
aggregate correlation at $0.43$, $0.38$ and $0.40$ respectively, each with a 95\%
CI, confirming the aggregate correlation is not driven by any subset of poems. Per-author
agreement is weaker: across the ten gold authors with at least three sampled poems the
agreement is $\rho{=}0.13$ (95\% CI $[-0.55,0.69]$, $n{=}10$), indistinguishable from
nothing at that sample size; we report it rather than search for the aggregation that
flatters the instrument. The gold set is confined to \emph{shi}; nothing here licenses
the analyser for \emph{ci} or prose.


\begin{thebibliography}{9}
\bibitem[Kim(2014)]{kim2014convolutional}
Y.~Kim.
\newblock Convolutional neural networks for sentence classification.
\newblock In \emph{EMNLP}, 2014.

\bibitem[Zhou et al.(2022)]{shiziren2022}
A.~Zhou, C.~Sang, Y.~Zhang, and M.~Lu.
\newblock 诗人密码: 唐诗作者身份识别 (The Poet Code: Tang Poetry Authorship
Identification).
\newblock \emph{Journal of Chinese Information Processing
(中文信息学报)}, 36(6):162--170, 2022. DOI: 10.3969/j.issn.1003-0077.2022.06.017.

\bibitem[Zhou et al.(2021)]{zhou2021lda}
A.~Zhou, Y.~Zhang, H.~Wei, and M.~Lu.
\newblock LDA-Transformer Model in Chinese Poetry Authorship Attribution.
\newblock In \emph{Information Retrieval: 27th China Conference, CCIR 2021,
Dalian, China, October 29--31, 2021, Proceedings (LNCS, vol.~13026)}, pages
59--73. Springer, 2021. DOI: 10.1007/978-3-030-88189-4\_5.

\bibitem[Fang(2011)]{fang2011cld}
A.~C.~Fang, W.-y.~Li, and J.~Cao.
\newblock In search of poetic discourse of classical Chinese poetry: an
  imagery-based stylistic analysis of Liu Yong and Su Shi.
\newblock \emph{Chinese Language and Discourse}, 2(2):232--249, 2011.
  John Benjamins. DOI: 10.1075/cld.2.2.04fan.
\end{thebibliography}

\end{document}
